\section{阶段拆分沿依赖图而不是功能清单}

里程碑之间存在硬依赖：没有可观察日志就难以调试磁盘，没有稳定装载就无法验证
模式切换，没有异常入口就不应扩大内核复杂度，没有内存所有权就不能安全创建
进程。阶段顺序由“下一机制需要哪些已证明状态”决定。

\begin{pdfdiagram}
  \node[os state] (observe) {复位与串口\\取得控制和观测};
  \node[os block,right=of observe] (load) {磁盘装载\\取得后续字节};
  \node[os state,right=of load] (long) {长模式与 ELF\\建立 C++ 入口};
  \node[os block,below=of long] (exception) {异常与内存\\建立保护和所有权};
  \node[os state,left=of exception] (device) {中断与驱动\\建立异步事件};
  \node[os block,left=of device] (user) {进程与文件\\建立用户环境};
  \draw[os arrow] (observe) -- (load);
  \draw[os arrow] (load) -- (long);
  \draw[os arrow] (long) -- (exception);
  \draw[os arrow] (exception) -- (device);
  \draw[os arrow] (device) -- (user);
\end{pdfdiagram}
\diagramnote{每个节点只增加一组新不变量；箭头表示下一阶段可以依赖的已验收能力。}

版本号表达教学能力而非市场发布。一个版本只有在实现、失败路径、测试、教材和
发布证据同时完成后才结束，不能以“正常路径在本机运行一次”作为迁移条件。

\section{需求使用可观察语言固定验收边界}

“实现磁盘驱动”无法直接验收。可验证需求写成：给定符合格式的镜像，固件在
有限预算内读取指定 LBA，校验内容并输出阶段标记；给定 ERR、忙超时、短镜像
和越界范围，分别拒绝并留下稳定错误码。

\begin{longtable}{L{0.19\textwidth}L{0.27\textwidth}L{0.40\textwidth}}
\toprule
\textbf{需求组成} & \textbf{记录内容} & \textbf{磁盘阶段示例}\\
\midrule
\endfirsthead
\toprule
\textbf{需求组成} & \textbf{记录内容} & \textbf{磁盘阶段示例}\\
\midrule
\endhead
前置状态 & 已验证硬件和软件条件 & UART 可用、低端 RAM 区间保留\\
动作 & 新机制执行的事务 & IDE PIO 读取一个扇区\\
成功结果 & 状态与外部标记 & 校验通过并跳到 Stage 1\\
失败结果 & 分类与有限性 & BSY 超时、ERR、越界各自停止\\
证据 & 独立可重复检查 & 单元、随机、QEMU 与镜像审计\\
\bottomrule
\end{longtable}

\section{架构决策记录保留被放弃方案}

ADR 包含背景、约束、候选方案、决定、后果和重新评估条件。只写最终选择会让
后来者重复已经分析过的路线，也无法判断条件变化后是否应该改选。

例如选择 IDE PIO 而不是 AHCI，是因为当前 QEMU PC 基线、实现规模和学习
目标；代价是性能低且依赖传统设备。进入多队列与 DMA 学习阶段后，原决定的
适用条件变化，可以新增 ADR 替代，而不修改历史记录。

\section{模块文档把接口与失败语义一起公开}

每个模块记录职责、输入输出类型、所有权、并发上下文、硬件来源、不变量和错误
枚举。函数声明只告诉编译器参数类型，模块文档补足调用时序和生命周期。

设备驱动还要记录寄存器表、访问宽度、状态图、超时来源与复位策略；内存模块
记录地址空间、页大小、权限和 TLB 要求；调度模块记录锁顺序和允许切换点。

\section{测试矩阵由机制维度与错误维度交叉形成}

按测试名称逐个增加容易遗漏系统性边界。先列机制维度，再列输入类别：

\begin{longtable}{L{0.20\textwidth}L{0.21\textwidth}L{0.21\textwidth}L{0.26\textwidth}}
\toprule
\textbf{机制} & \textbf{正常} & \textbf{边界} & \textbf{故障}\\
\midrule
\endfirsthead
\toprule
\textbf{机制} & \textbf{正常} & \textbf{边界} & \textbf{故障}\\
\midrule
\endhead
地址区间 & 包含与相邻 & 空区间、最高地址 & 加法溢出、重叠保留区\\
ROM & 正确入口 & 精确末端 16 字节 & 错位移、错误尺寸\\
UART & 多字节发送 & 空字符串、最大预算 & 永不就绪、线路错误\\
ELF & 多个合法段 & 零 BSS、最大表项 & 截断、重叠、权限非法\\
页分配 & 分配释放 & 首尾页、耗尽 & 重复释放、保留页请求\\
\bottomrule
\end{longtable}

矩阵不是要求每格一个独立测试文件，而是确保性质与失败类别都有观察点。随机
测试适合跨大量组合，固定样例适合保留具有解释价值的边界。

\section{调试记录按假设与证据推进}

调试日志记录复现命令、最后已知阶段、候选原因、观测手段、结果和结论。一次
只改变一个假设，避免同时修改链接脚本、汇编和 QEMU 参数后无法知道真正原因。

裸机无输出时按外层向内检查：宿主进程是否启动、镜像尺寸与路径、复位固定
字节、入口反汇编、CPU 模式和段状态、栈地址、端口事务。每一步排除一类故障，
不以“多加几条日志”替代因果定位。

\section{发布记录把提交、工具与证据冻结}

每个版本记录提交哈希、工具版本、镜像哈希、串口标记、测试数量、已知限制和
下一阶段输入。发布不是复制一个二进制，而是保存重建它所需的环境与解释。

工具升级单独提交并执行全量回归。若 NASM 或 LLD 改变编码与布局，镜像哈希
变化可以追溯到工具决策，而不会与功能修改混杂。

\section{代码量统计只描述规模切片}

当前 231 行有效目标代码是按固定扩展名和注释口径得到的切片。它适合观察真实
系统代码随阶段增长，不适合比较语言生产率或判断质量。链接脚本虽不计入执行
代码，却仍是关键工程资产。

更有意义的进度指标包括已验证硬件事务、明确错误类别、性质测试数量、最大无界
等待数、未解析运行时符号和文档失效项。任何指标都可能被追逐而失真，所以最终
仍由机制级验收判定。

\section{代码评审沿数据流与所有权展开}

评审先确认外部输入来自何处、采用什么宽度、在哪一步验证，再沿地址计算、状态
转换和资源所有权追踪到输出。看到硬件访问就核对规范位，看到加法乘法就核对
溢出，看到指针就核对地址空间与生命周期，看到等待就核对终止条件。

风格检查属于基础层：类成员显式 \code{this->}，常量有项目与模块前缀，头源
分离，中文注释解释原因，固定宽度类型匹配协议。风格一致使评审注意力集中到
系统不变量。

\section{教材与代码采用同源引用和自动检查}

源码清单直接从工程路径输入，避免手工复制后失效；真实代码统计在教材构建前
生成；输入图检查防止章节漏接；PDF 构建检查链接、字体和版式。网站和手机
导出都使用同一 PDF 产物。

代码尚未实现的章节明确使用未来时态与设计约束，已实现章节给出当前字节、命令
和测试证据。随着 v0.2 开发，磁盘章节将加入真实源码清单和执行日志，不能让
设计说明长期冒充完成状态。

\section{一次阶段开发的完整闭环}

以 v0.2 为例，开发顺序固定为：

\begin{enumerate}[label=\arabic*.]
  \item 从 ATA 与 QEMU 机器模型提取端口、状态和时序约束。
  \item 写 ADR，确定 PIO、LBA 范围、Stage 1 格式与加载地址。
  \item 为纯范围计算和头解析先写宿主单元与随机性质测试。
  \item 实现有界设备状态机和稳定错误枚举。
  \item 构造成功、短镜像、错误状态和忙超时的 QEMU 场景。
  \item 审计最终 ROM 与磁盘布局，保存串口证据。
  \item 更新模块、调试、教材与发布文档。
  \item 在干净环境执行完整回归，达到退出标准后标记版本。
\end{enumerate}

闭环使每个新机制同时拥有知识背景、代码实现、失败行为和自动证据。项目规模
增长后仍重复同一节奏，避免“先堆功能，最后补文档和测试”形成无法偿还的系统
债务。
